Model dokumentowy dla MongoDB stanowi denormalizowane przekształcenie modelu
relacyjnego, dostosowane do typowych wzorców dostępu w aplikacjach
streamingowych. Schemat logiczny obejmuje sześć kolekcji głównych --
\texttt{users}, \texttt{content}, \texttt{watch\_history}, \texttt{ratings},
\texttt{payments} oraz \texttt{my\_list} -- w których część encji powiązanych
relacjami jeden-do-wielu została zagnieżdżona w postaci poddokumentów. Schemat
graficzny przedstawiono na rysunku~\ref{fig:schema_mongo}.

\begin{figure}[H]
    \centering
    % \includegraphics[width=1.0\textwidth]{schemas/schema_mongo.png}
    \caption{Schemat dokumentowy dla MongoDB}
    \label{fig:schema_mongo}
\end{figure}

Kolekcja \texttt{users} agreguje w pojedynczym dokumencie informacje o koncie
użytkownika, tablicę profili (\texttt{profiles[]}) oraz aktualny obiekt
subskrypcji (\texttt{subscription\{\}}). Decyzja o zagnieżdżeniu profili
wynika z faktu, że ich liczba w obrębie konta jest ograniczona (najczęściej
od dwóch do pięciu) oraz że są one praktycznie zawsze pobierane wspólnie
z danymi użytkownika podczas operacji logowania lub przełączania profilu.
Subskrypcja przechowywana jest jako pojedynczy zagnieżdżony obiekt
reprezentujący stan bieżący, podczas gdy historia płatności została
wydzielona do osobnej kolekcji \texttt{payments} z uwagi na rosnący
w czasie wolumen rekordów oraz brak konieczności pobierania ich razem
z danymi użytkownika.

Kolekcja \texttt{content} skupia w pojedynczym dokumencie kompletny opis
filmu lub serialu wraz z zagnieżdżoną tablicą obsady (\texttt{cast[]}) oraz,
w przypadku seriali, dwupoziomową hierarchią sezonów i odcinków
(\texttt{seasons[]} zawierającą \texttt{episodes[]}). Tablica obsady
przechowuje denormalizowaną kopię podstawowych danych każdej osoby --
imię, nazwisko, rolę oraz, dla aktorów, nazwę postaci i pozycję w napisach --
co eliminuje konieczność wykonywania operacji \texttt{\$lookup} przy
typowym scenariuszu pobrania szczegółów tytułu. Gatunki przechowywane są
jako tablica wartości tekstowych (\texttt{genres[]}), co umożliwia
wykorzystanie indeksu wielowartościowego (\textit{multikey index}) do
efektywnego wyszukiwania treści po wybranej kategorii. Pole \texttt{metadata},
typu \texttt{object}, agreguje dodatkowe atrybuty marketingowe i techniczne,
omówione szczegółowo w kolejnym podrozdziale.

Kolekcje rejestrujące aktywność użytkownika -- \texttt{watch\_history},
\texttt{ratings} oraz \texttt{my\_list} -- zostały utrzymane jako odrębne
struktury, mimo że w modelu hierarchicznym stanowiłyby naturalne kandydatury
do zagnieżdżenia w obrębie dokumentu profilu. Decyzja o ich wydzieleniu
wynika z dwóch przesłanek: rosnącego w czasie wolumenu rekordów, który
w przypadku zagnieżdżenia mógłby przekroczyć limit szesnastu megabajtów
przypadający na pojedynczy dokument BSON, oraz konieczności efektywnego
wyszukiwania po identyfikatorze treści, na przykład w celu wyznaczenia
średniej oceny lub liczby unikalnych widzów dla danego tytułu. Powiązania
z dokumentami nadrzędnymi realizowane są przez referencje
(\texttt{profile\_id}, \texttt{content\_id}), z indeksami złożonymi
zaprojektowanymi pod konkretne wzorce zapytań -- typowo
\texttt{(profile\_id, started\_at)} dla historii oglądania oraz
\texttt{(content\_id, score)} dla ocen.

Kolekcja \texttt{payments} została wydzielona z dokumentu użytkownika,
ponieważ liczba transakcji per subskrypcja może być znacząca, a dane te
są przetwarzane głównie w zapytaniach analitycznych (raporty przychodów,
analiza skuteczności metod płatności), niezwiązanych bezpośrednio
z operacjami na koncie użytkownika. Powiązanie z subskrypcją realizowane
jest przez referencję \texttt{subscription\_id}, mapującą się na identyfikator
zagnieżdżony w dokumencie \texttt{users}.

Pole \texttt{\_id} stanowi domyślny identyfikator każdego dokumentu, na
którym automatycznie zakładany jest unikalny indeks. Dla zachowania
kompatybilności semantycznej z modelem relacyjnym zachowano typ
\texttt{bigserial} (mapowany na 64-bitową liczbę całkowitą), zamiast
domyślnego typu \texttt{ObjectId}, co ułatwia generowanie identyfikatorów
zewnętrznym skryptem oraz zapewnia bezpośrednią porównywalność z modelami
SQL podczas testów wydajnościowych.
